% =============================================
% INTRODUCTION GÉNÉRALE
% =============================================
\chapter*{Introduction Générale}
\addcontentsline{toc}{chapter}{Introduction Générale}
\thispagestyle{fancy}

\section*{Contexte général}

Dans un monde où le développement logiciel s'accélère et où les organisations gèrent
simultanément de multiples projets impliquant des équipes distribuées, la maîtrise des
flux de travail et de la collaboration devient un enjeu stratégique majeur. Les outils
de gestion de projets existants, bien que variés, présentent souvent des limites en
termes de personnalisation, d'intégration avec des systèmes d'authentification d'entreprise
et de gestion fine des rôles et des responsabilités.

C'est dans ce contexte que s'inscrit le présent projet de fin d'année, réalisé
à l'\textbf{École Marocaine des Sciences de l'Ingénieur (EMSI)}, visant à concevoir
et développer une solution back-end complète pour la \textbf{gestion de projets collaboratifs}.

\section*{Problématique}

Comment concevoir une API RESTful robuste, sécurisée et extensible permettant
de gérer l'ensemble du cycle de vie d'un projet logiciel, tout en supportant
une hiérarchie de rôles claire et en garantissant la conformité aux accords
de niveau de service (SLA) définis ?

\section*{Objectifs du projet}

Ce projet vise à atteindre les objectifs suivants :

\begin{itemize}[leftmargin=*, label=\textcolor{emsiRed}{\ding{118}}]
  \item Concevoir une architecture d'API REST maintenable et évolutive avec \textbf{Laravel}.
  \item Implémenter un système de \textbf{gestion multi-rôles} (Manager, Chef de Projet, Développeur).
  \item Intégrer une authentification centralisée et sécurisée via \textbf{Keycloak/JWT}.
  \item Gérer le cycle de vie complet des projets et des tâches avec suivi des \textbf{SLA}.
  \item Fournir des \textbf{statistiques et exports} de données pour le pilotage décisionnel.
  \item Automatiser la création d'espaces de travail de projet (\textit{workspace scaffolding}).
\end{itemize}

\section*{Structure du rapport}

Ce rapport est organisé en cinq chapitres :

\begin{enumerate}[leftmargin=*, label=\textcolor{emsiRed}{\textbf{Chapitre \arabic*.}}]
  \item \textbf{Contexte et Présentation du Projet} --- Présentation de l'organisme d'accueil,
    de l'environnement technique et de la méthodologie adoptée.
  \item \textbf{Analyse et Spécification des Besoins} --- Identification des acteurs,
    des cas d'utilisation et des exigences fonctionnelles et non-fonctionnelles.
  \item \textbf{Conception} --- Modélisation UML, architecture logicielle et modèle de données.
  \item \textbf{Réalisation et Tests} --- Présentation de l'implémentation, des fonctionnalités
    clés et des tests effectués.
  \item \textbf{Conclusion et Perspectives} --- Bilan du travail réalisé et pistes d'évolution futures.
\end{enumerate}
